\chapter{Introducción}
\label{chap:introduccion}

Las criptomonedas han revolucionado el panorama financiero, ofreciendo nuevas oportunidades de inversión y especulación a corto, mediano y largo plazo \parencite{baur_bitcoin_2017}. Sin embargo, la volatilidad inherente a estos activos y la complejidad de los mercados digitales presentan desafíos significativos para los inversores \parencite{katsiampa_volatility_2017}. A diferencia de los mercados bursátiles tradicionales, los mercados de criptomonedas operan de forma ininterrumpida, los siete días de la semana, a escala global y sin un regulador centralizado. Esta naturaleza descentralizada, sumada a la velocidad con la que los precios pueden variar en cuestión de minutos, hace que la toma de decisiones humana por sí sola sea insuficiente para competir con los sistemas automatizados que dominan actualmente estos mercados. Mientras que en la bolsa tradicional un inversor minorista podía analizar la situación con calma durante el horario de mercado, en el ecosistema cripto esta ventana se ha cerrado: el mercado no duerme, las oportunidades son efímeras y los errores de juicio tienen consecuencias inmediatas. A esto se añade que la barrera de entrada se ha reducido drásticamente gracias a plataformas de fácil acceso, lo que ha atraído a millones de inversores sin experiencia previa que se enfrentan a un entorno hostil sin las herramientas ni los conocimientos necesarios para navegarlo de forma racional.

La necesidad de este proyecto radica, en primer lugar, en la dificultad inherente de los usuarios para reaccionar en tiempo real ante estas rápidas fluctuaciones del mercado. A esto se suma una carencia generalizada de educación financiera básica, especialmente entre los jóvenes y los inversores minoristas \parencite{lusardi_economic_2014}. En la actualidad, resulta complejo iniciarse en el mundo de la inversión de forma segura y sin arriesgar capital real. Este problema se ve agravado por la sobreabundancia de ruido mediático y atención pública desmedida en torno a estos activos. Como demuestra la literatura reciente, esto fomenta expectativas irreales que empujan a los usuarios inexpertos (especialmente a los más jóvenes) a operar de forma altamente especulativa y con excesiva frecuencia, lo que habitualmente desemboca en menores rendimientos y la pérdida paulatina de sus ahorros \parencite{hasso_who_2019}. Por otro lado, el trading algorítmico ofrece la gran ventaja de eliminar el sesgo emocional de las decisiones de inversión, aportando una mayor disciplina y consistencia matemática en la ejecución de las estrategias \parencite{chan_algorithmic_2013}. Sin embargo, la alta barrera de entrada, derivada de la profunda complejidad técnica y la falta de acceso a herramientas de simulación adecuadas, dificulta enormemente el aprendizaje \parencite{treleaven_algorithmic_2013}.

En este contexto, las \textit{estrategias de trading} constituyen la respuesta natural a estos problemas. Una estrategia no es más que un conjunto de reglas objetivas y reproducibles basadas en indicadores técnicos como la media móvil, el RSI o el MACD— que determinan de forma autónoma cuándo comprar y cuándo vender un activo, eliminando así la subjetividad y la reacción emocional del inversor humano \parencite{murphy_technical_1999}. El valor real de las estrategias, sin embargo, no reside únicamente en su ejecución en tiempo real, sino en la posibilidad de validarlas previamente sobre datos históricos mediante el \textit{backtesting}: un proceso que permite medir su rendimiento pasado antes de comprometer capital real. Esta capacidad de experimentación controlada resulta fundamental para el aprendizaje, pues permite al inversor iterar sobre sus hipótesis, detectar sus errores y desarrollar un criterio crítico sin consecuencias económicas. A pesar de ello, el diseño, la implementación y la validación de estas estrategias sigue siendo una tarea reservada a perfiles con conocimientos avanzados de programación y matemáticas financieras, lo que excluye a la inmensa mayoría de los inversores interesados en adoptarlas.

Para abordar esta problemática, este proyecto se centra en el desarrollo de una plataforma de trading algorítmico orientada a reducir esta brecha técnica, democratizando el acceso al trading cuantitativo. El sistema proporciona a los estudiantes un entorno de simulación seguro (\textit{Paper Trading}) donde pueden experimentar, aprender de sus errores y validar sus hipótesis sin riesgo financiero. Además, la plataforma centraliza la exportación de resultados en formato CSV para su auditoría. Asimismo, para aquellos con un perfil más técnico, el sistema proporciona una base sólida para la recopilación de datos históricos de múltiples monedas y la aplicación práctica de modelos de Inteligencia Artificial, ya sea entrenando algoritmos con indicadores básicos o retroalimentándolos con los resultados de las estrategias del usuario.

Finalmente, para facilitar la curva de aprendizaje, el formato de las estrategias se basa en el estándar utilizado por \textit{Freqtrade} \parencite{freqtrade_development_team_freqtrade_2023}, una plataforma de código abierto escrita en Python. Este formato permite a los usuarios definir sus propias reglas operativas de forma intuitiva. Las estrategias se implementan como scripts que interactúan con el motor de cálculo para recibir datos de mercado, procesar señales y enviar órdenes, permitiendo al estudiante centrarse puramente en la lógica financiera mientras el sistema gestiona la complejidad subyacente.
